% AUTOMATISCH ERZEUGT aus dem SysML-Modell mit tools/generate_docs.py.
% Nicht von Hand aendern: Inhalte im Modell pflegen, Layout im Kapitel.
REQ-TRN-001 & Nicht MVP & Materialumschlag wird im Zielsystem modelliert, aber im ersten MVP noch nicht voll implementiert. \\
REQ-MAP-001 & MVP & Ohne Grundkarte ist der geplante Navigationskern nicht nachweisbar. \\
REQ-MAP-002 & MVP & Laufzeitaktualisierung der Umgebung ist Kern der Projektidee. \\
REQ-LOC-001 & MVP & Lokalisierung ist Voraussetzung fuer zielgerichtete Navigation. \\
REQ-NAV-001 & MVP & Autonomes Zielanfahren ist zentraler Funktionsnachweis. \\
REQ-NAV-002 & MVP & Hindernisbehandlung prueft die Kopplung von Wahrnehmung und Navigation. \\
REQ-NAV-004 & MVP & Mehrere Raeume und Flur bilden den Demonstrationskern. \\
REQ-NAV-005 & MVP & Offene Tueren sind zentrale Randbedingung der Umgebung. \\
REQ-SAFE-001 & MVP reduziert & Fuer MVP mindestens sicheres Stoppen; spaetere Safety-Auslegung bleibt offen. \\
REQ-UI-001 & MVP & Die Web-App ist die erste reale Auftragsschnittstelle. \\
REQ-UI-002 & MVP reduziert & Status, betrieblicher Stopp mit manueller Fortsetzung und Rueckkehr zur Ladestation werden als Grundbefehle betrachtet. \\
REQ-PWR-001 & MVP reduziert & Ohne Ladezustand ist der 2-h-Nachweis aus REQ-PWR-002 nicht moeglich; im MVP genuegt Erkennen und Melden eines niedrigen Ladezustands, ein genaues Energiemodell folgt spaeter. \\
REQ-PWR-002 & MVP & 2 h Betriebsdauer decken einen 90-min-Unterrichtsblock mit Reserve ab. \\
REQ-PERF-001 & MVP & Niedrige Geschwindigkeit reduziert Risiko und erleichtert Unterricht und Tests. \\
REQ-STA-001 & MVP & Statusmeldung ist fuer Testbarkeit und Bedienbarkeit erforderlich. \\
